\section{Logical reductions, rejected deletions, and proof strength}
\label{sec:logical-audit}

This section separates three questions: equivalence of abstract proof
calculi, the number of states obtained by compiling them, and formalization
of the resulting arithmetic halting equivalence.  A reduction in one of
these measures does not settle the other two.

\subsection{An equivalent split of the first quantifier head}

Write $f=1$ for the distinguished contradiction code and put
\[
 X=\mathsf U\bigl(D,\mathsf U(D,H)\to f\bigr).
\]
The first PA196 head is
\[
 O(D,H)=(X\to f)\to H.
\]
Consider replacing it by the two families
\[
 B(D,H)=(H\to f)\to X,
 \qquad E(A)=f\to A.
 \tag{Split}
\]
All remaining heads, the guarded gate, modus ponens, and closed
generalization are unchanged.

\paragraph{Split-head theorem.}
The original calculus and the calculus with \textup{(Split)} derive exactly
the same numeric formulas.  The conversions are uniform finite
implicational proof templates and hence induce primitive-recursive
transformations of coded derivations.

\paragraph{Proof.}
The restricted $\rho_8$ family supplies $K$, $S$, and Peirce's law without
using the first quantifier head or explosion.  This is explicit in
\texttt{M6/Logic/RestrictedImplication.lean}: its abstract \texttt{Rules}
structure contains only modus ponens, the restricted row, and the three
code-order conditions.  The PA196 interpretation of those rules uses only
its $\rho_8$ constructor.  In particular, ordinary implicational deduction
and
\[
 ((P\to Q)\to P)\to P
\]
are available before either side of the proposed conversion.

In the new calculus assume $p:X\to f$.  Under the further assumption
$n:H\to f$, the new head $B(D,H)$ gives $X$, then $p$ gives $f$, and
$E(H)$ gives $H$.  Discharge $n$ to obtain $(H\to f)\to H$ and apply
Peirce's law with $(P,Q)=(H,f)$.  Discharging $p$ now gives $O(D,H)$.

Conversely, the old first head supplies explosion.  For any $A$, take
\[
 Y=\mathsf U\bigl(0,\mathsf U(0,A)\to f\bigr),
 \qquad (Y\to f)\to A.
\]
From the assumption $f$, $K$ gives $Y\to f$; the displayed head gives $A$.
Discharge $f$ to obtain $E(A)$.  To obtain $B(D,H)$, assume $n:H\to f$ and,
inside this context, $q:X\to f$.  The old head gives $H$, then $n$ gives
$f$, and the already derived $E(X)$ gives $X$.  Discharge $q$, apply
Peirce's law with $(P,Q)=(X,f)$, and discharge $n$.  The result is
$(H\to f)\to X$, as required.  No quantifier rule occurs in either
conversion.\hfill$\square$

The accompanying \texttt{check\_split\_head.py} constructs exact Hilbert
proof DAGs and checks every node using $K$, $S$, Peirce, assumptions, and
modus ponens.  The new-to-old and old-to-new templates each have 77 nodes;
the old-to-explosion template has 17 nodes.  These counts measure the
particular checked certificates, not minimal proof lengths.  Their status
is a finite syntactic verification over the displayed derived
implicational basis, not a full new Lean build of the tape machine.

Keeping the explicit explosion family is essential.  The contraposed
head $B(D,H)$ alone is not a replacement justified by classical
contraposition: that argument already assumes that code $1$ acts as
falsity.  Indeed, interpret implication materially, every membership atom as true,
every equality atom as false, and every quantifier as the identity
operation on truth values.  Code $0$ is true, code $1$ is true, and code
$3$ is false.  Every retained head and the guarded gate, together with
$B(D,H)$, is valid under this evaluation; generalization preserves
validity.  But $1\to3$ is false.  The proof above identifies the exact
additional family that repairs this obstruction.

\subsubsection{The corresponding command words}

The split adds one head.  An axiom command consequently pushes ten items:
the saved top, one generalization, $\rho_8$, and seven heads.  The guarded
gate remains the new top.  If $S$ denotes Shift, the corrected generic
stack words are
\[
\begin{aligned}
 \mathrm{pop} &:\quad \mathsf{AX}(0,0,0)\,S^{11},\\
 \mathrm{push}&:\quad \mathsf{AX}(0,0,0)\,S^{9}\,\mathsf{MP}(0),\\
 \mathrm{discard}&:\quad
 \mathsf{MP}(1)\,\mathsf{AX}(1,0,0)\,S^{10}\,\mathsf{MP}(0).
\end{aligned}
\]
The generalization snapshot identity becomes
\[
 \mathsf{AX}(c,x,y)\,S^{10}(t,w,d,h)
   =\bigl(\mathsf U(d,h),w,0,t\bigr).
\]
For a head numbered $j$ from the top among the eight nongeneralization
heads, $1\le j\le8$, the generic selector is: discard, execute
$\mathsf{AX}$ with the desired arguments, apply $S^{j+1}$, and pop
$9-j$ times.  The gate selector discards, executes $\mathsf{AX}$, and
pops ten times.  These are the PA196 stack identities with the push count
increased by one; the position of each formula must also follow the
chosen emission order.  The two new head families have implication as
their outer constructor, so their codes are $0$ modulo $4$ and cannot
equal the halt code $1$.  Taking a clean store $d=h=0$ makes the parameters
of $B(D,H)$ and $E(A)$ unrestricted.

These command-word adjustments are necessary even when the emitted
formula families are logically equivalent.  A candidate table that
changes the output count but retains the old selectors has not thereby
inherited the original certificate compiler.

\subsection{Why neither membership rewrite can simply be deleted}

The two direct rewrite families are
\[
\begin{aligned}
 L(a,b,c)&:\quad
 \mathsf E(a,b)\to\bigl(\mathsf M(a,c)\to\mathsf M(b,c)\bigr),\\
 R(a,b,c)&:\quad
 \mathsf E(a,b)\to\bigl(\mathsf M(c,a)\to\mathsf M(c,b)\bigr).
\end{aligned}
\]
They correspond to PA196 heads 3 and 4, with the harmless antecedent
exchange used in head 3.  The following countermodels rule out the
proposal to delete either one and derive it from all the remaining
numeric axiom families.

Let $\operatorname{Bit}(s,k)$ mean that bit $k$ of $s\in\mathbb N$ is one.
Let $f_1$ exchange $0$ and $1$, and let every other $f_i$ be the identity.
Thus every $f_i$ is a finite-support permutation.  Variables are assigned
natural-number values $x_i$; implication and quantifiers have their
ordinary classical meanings, and
$\mathsf E(a,b)$ means $x_a=x_b$.  Numeric code $0$ denotes truth, as in
the usual decoder.

\paragraph{Delete $L$.}
Interpret membership by
\[
 \mathsf M(a,b)\quad\Longleftrightarrow\quad
 \operatorname{Bit}\bigl(f_a(x_a),x_b\bigr).
 \tag{L-model}
\]
The family $R$ is valid because equal values in its second membership
coordinate can be substituted.  The quantifier heads, existential
equality witness, $\rho_8$, and both atomic vacuity families are valid.
Moreover, code $1$ still denotes
$\operatorname{Bit}(x_0,x_0)$, which is always false.

The guarded gate remains valid for every formula instance, including the
case where its two ordinary labels coincide.  When $a\ne b$, the values
of the bound variable satisfying $\mathsf M(a,b)$ form the finite bitset
of $f_a(x_a)$.  When $a=b$, they form
\[
 \{n:\operatorname{Bit}(f_a(n),n)\}.
\]
This set is finite: outside the finite support of $f_a$, the condition is
$\operatorname{Bit}(n,n)$ and is false.  Add the distinguished point
$x_b$ to this finite candidate set and retain precisely the elements
satisfying the gate's arbitrary formula.  If $s$ codes the selected
finite set, assign $x_c=f_c^{-1}(s)$.  The guard $c\ne a,b$ and the
self-freshness of numeric formula code $c$ ensure that this assignment
changes neither the old parameters nor the selected property.  It
therefore supplies exactly the required pointed separator.  Equality
and membership vacuity under the same guard are immediate.

Nevertheless, assign $x_1=x_2=x_3=0$.  Then $\mathsf E(1,2)$ and
$\mathsf M(1,3)$ hold while $\mathsf M(2,3)$ fails.  Thus $L(1,2,3)$ is
false and cannot be derived from the retained families.

\paragraph{Delete $R$.}
Instead interpret membership by
\[
 \mathsf M(a,b)\quad\Longleftrightarrow\quad
 \operatorname{Bit}\bigl(x_a,f_b(x_b)\bigr).
 \tag{R-model}
\]
The family $L$ is now valid.  All other retained heads, both atomic
vacuities, and the false interpretation of code $1$ remain as above.
For $a\ne b$, the gate's candidate domain is the inverse image under
$f_b$ of the finite bitset of $x_a$.  For $a=b$, it is the finite set
\[
 \{n:\operatorname{Bit}(n,f_b(n))\}.
\]
After adding the distinguished point and filtering by the arbitrary
property, assign to $x_c$ the bitset of the images of the selected values
under $f_b$.  This proves the entire gate schema in this model.

Set $x_1=x_2=0$ and $x_3=2$.  Then $\mathsf E(1,2)$ and
$\mathsf M(3,1)$ hold, since bit $1$ of $2$ is one, while
$\mathsf M(3,2)$ fails, since bit $0$ of $2$ is zero.  Hence $R(1,2,3)$
is not derivable from its retained families.

These are generalized semantics for the actual numeric calculus, not
ordinary structures with one uniform binary predicate.  That distinction
is deliberate: the calculus's primitive proof rules are modus ponens
and closed generalization.  An unrestricted substitution rule for
arbitrary numeric formula codes is not available.  Since every retained
axiom family is valid and the two primitive rules preserve validity,
these semantics suffice to prove nonderivability.  They exclude a
\emph{derived-head justification} for either deletion.  They do not
exclude an entirely new interpretation proving that a smaller calculus
has the same consistency strength.

The script \texttt{check\_rewrite\_countermodels.py} verifies the two
explicit countervaluations and finite controls of the constructions.
For labels $0,\ldots,4$ and parameter values $0,\ldots,7$, it checks
8,000 retained-rewrite cases and 23,056 separator cases for each model,
exhausting every property restriction on each relevant finite candidate
set.  This is a check of the stated finite ranges.  The full-schema
argument is the mathematical finite-support construction above.

\subsection{What the iPRA work can contribute}

The native scalar core in the supplied \texttt{ConciseIPRA} project~\cite{ipra} is
bi-interpretable with intuitionistic primitive recursive arithmetic.
Using that core unchanged would target a different consistency
statement.  Its most direct useful modification is precise: retain its
native terms, formula constructors, logical rules, and freshness
conditions, but extend its progress rule from literal equations to
arbitrary native formulas:
\[
 \frac{\Gamma,B(z)\vdash B(Sz)}
      {\Gamma\vdash Q_y B(y)},
 \qquad
 Q_y B(y)\equiv B(0)\to\forall y\,B(y),
 \qquad z\text{ fresh}.
 \tag{Full progress}
\]
The ordinary realization of this extended core is sound in Heyting
arithmetic, using full formula induction.  Conversely, the existing
native definitions of the logical connectives and primitive-recursive
functions, together with \textup{(Full progress)} followed by closure
elimination, supply induction for every translated Heyting-arithmetic
formula.  This gives a natural HA-strength source calculus.  The usual
negative translation from PA to HA and the inclusion of intuitionistic
proofs into classical proofs then supply the intended consistency
comparison, once implemented on the selected proof codes.

This is a concrete high-level avenue, not a delivered smaller tape
machine.  A compilation must still count the implementation of terms,
substitution, freshness, contexts, certificate enumeration, and
blank-tape startup.  Python token counts and numbers of logical rule
families do not determine binary single-tape state counts.

\subsection{The reverse internal PA implication}

Let $H_{196}$ be the exact arithmetic halting sentence for the literal
PA196 table, and let $N_{196}=\neg H_{196}$.
The upstream declaration
\texttt{PA196.Arithmetic.pa\_proves\_nonhalting\_consistency}
proves
\[
 \mathrm{PA}\vdash N_{196}\to\operatorname{Con}(\mathrm{PA}).
\]
The source's published theorem inventory does not include the converse.
Its external nonhalting proof uses a hereditarily finite-list model in
Lean's metatheory.  That external soundness proof is not itself a PA
proof of the reverse implication.

There is a direct syntactic route to the missing direction.  Translate
numeric membership to $\operatorname{Bit}(x_a,x_b)$, numeric equality to
arithmetic equality, and decode the transposed quantifier constructor
correctly.  For every arithmetic formula $\varphi(x)$, PA proves
\[
 \forall a,b\;\exists s\;\forall x\quad
 \operatorname{Bit}(s,x)\leftrightarrow
 \bigl((\operatorname{Bit}(a,x)\lor x=b)\land\varphi(x)\bigr).
 \tag{PC}
\]
For a detailed template, set $m=\max(a,b+1)$ and use full induction on
$k\le m$ to construct a binary sum agreeing with the desired property
on the first $k$ bits.  At stage $k$, excluded middle for
$\varphi(k)$ determines whether to add $2^k$.  The implication
$\operatorname{Bit}(a,x)\Rightarrow x<a$ bounds every possible selected
value.  The formula $\varphi$ is inserted syntactically into this one
induction proof template.  Consequently the map from $\varphi$ to its
PA proof of \textup{(PC)} is primitive recursive; it does not evaluate
$\varphi$ or introduce a uniform arithmetic truth predicate.

Pointed comprehension supplies PA proof templates for the guarded
gate.  The other heads have fixed logical proof templates.  A
primitive-recursive induction on a numeric derivation then translates
any derivation of code $1$ into a PA refutation.  The command interpreter
can carry these proof codes in its register and stack invariant.  At a
halt test $t+h=1$, one of the two proved codes is $1$, so a refutation is
extracted.  This replaces semantic truth by syntactic proof objects
throughout the invariant.

The physical step is an internal version of the existing external
\texttt{blank\_halting\_iff\_source} theorem.  Positive forward macro
simulation is sufficient: if the literal machine halts by time $n$
and the source does not halt in its first $n+1$ steps, concatenate
$n+1$ positive macro simulations to obtain a literal execution longer
than $n$, contradicting deterministic halting.  This uses bounded
induction and the macro contracts, not an unbounded termination
argument.  A completed implementation of these constructions would
prove
\[
 \mathrm{PA}\vdash H_{196}\to
 \operatorname{Pr}_{\mathrm{PA}}(\ulcorner\bot\urcorner),
 \quad\text{hence}\quad
 \mathrm{PA}\vdash\operatorname{Con}(\mathrm{PA})\to N_{196}.
\]
This paragraph describes the proof construction and identifies its
formalization tasks.  It is not a claim that the missing arithmetic Lean
theorem has been supplied or checked in this work.

\subsection{A stronger obstruction to replacing the first head by reflexivity}

Even adding explicit explosion does not justify replacing the old first
head by the shorter family
\[
 \mathsf U(D,H)\to H.
 \tag{T}
\]
Here is a countermodel to that proposed derived-head argument.  A world
is a natural-number assignment together with a stage $s\in\{0,1\}$.
The quantifier $\mathsf U(i,\varphi)$ ranges over all worlds obtained by
changing $x_i$ arbitrarily and replacing $s$ by any $s'\ge s$.
Implication is material at each world.  Equality is ordinary numerical
equality.  At stage $0$ every membership atom is false; at stage $1$,
$\mathsf M(a,b)$ has its usual bit interpretation
$\operatorname{Bit}(x_a,x_b)$.  Code $1$ is false at both stages.

The accessibility relation for each quantifier is reflexive, so
\textup{(T)} holds.  Relations for different labels commute, since their
two coordinate changes and their final nondecreasing stage can be chosen
in either order.  Distribution holds for universal modalities.  The
equality witness holds by choosing the quantified value equal to the
other value.  Both membership rewrite families hold: at stage $0$ their
membership antecedents are false, and at stage $1$ they are ordinary
equality substitutions.  Atomic vacuity at a distinct label also holds:
equality is unchanged; a true membership atom already lies at stage $1$
and remains unchanged.  Every explosion instance holds because code
$1$ is false.

The guarded gate is valid at both stages.  Its equality and membership
vacuities were just checked.  To refute its universal $Q$ clause, choose
the separating variable at stage $1$.  All its subsequent quantifiers
then remain at stage $1$, where ordinary finite pointed separation
supplies a witness.  The formula-code freshness property has its usual
meaning at that stage.  The separating set makes both $S$ and $P$ true,
so $Q$ is false and the whole gate is true.

But take a stage-$0$ world with $x_2=1$, $x_3=0$, put
$H=\mathsf M(2,3)$, and take $D=4$.  Here $H$ is false.  After advancing
to stage $1$, $\mathsf U(4,H)$ is true, since changing $x_4$ leaves the
true assertion $\operatorname{Bit}(1,0)$ unchanged.  Hence
$\exists x_4\,\mathsf U(4,H)$ is true at the initial world, and the old
first head is false there.  Modus ponens and closed generalization
preserve validity in this semantics.  Consequently the first head is
not derivable from \textup{(T)}, explosion, and the remaining numeric
axiom families.  This is a stronger obstruction than merely observing
that a proposed replacement fails to make code $1$ behave as falsity.
